% Theory. Every statement below is either cited (Barber et al. 2023; Gibbs and Candes 2021) or proved in full.

\section{Theory}\label{sec:theory}

We use the notation of Section~\ref{sec:method}: weights $p_{t,s}$ that depend on the features $E_{1:t}$ only, $m_t=\max_{s<t}p_{t,s}$ and $n_{\mathrm{eff},t}=1/\sum_s p_{t,s}^2$.

\paragraph{The threshold.}
For a miscoverage level $\alpha_t$ the threshold is the weighted quantile of the past scores at level
\[
 \ell_t=(1-\alpha_t)(1+m_t),\qquad
 q_t=\inf\Bigl\{x:\ \textstyle\sum_{s<t}p_{t,s}\mathbf 1\{R_s\le x\}\ge \ell_t\Bigr\},
\]
with the conventions $q_t=+\infty$ if $\ell_t>1$ and $q_t=-\infty$ if $\ell_t\le 0$.
The factor $1+m_t$ is the finite-sample correction of Barber et al.\ (2023): it equals the
quantile, at level $1-\alpha_t$, of $\sum_{s<t}\tilde w_{t,s}\delta_{R_s}+\tilde w_{t,t}\delta_{+\infty}$ with
$\tilde w_{t,s}=p_{t,s}/(1+m_t)$ and $\tilde w_{t,t}=m_t/(1+m_t)$. For uniform weights on $n$ scores
this is the usual split-conformal level $(1-\alpha)(1+1/n)$.
The level follows the online rule $\alpha_{t+1}=\alpha_t+\gamma(\alpha-\mathrm{err}_t)$,
$\mathrm{err}_t=\mathbf 1\{R_t>q_t\}$, $\alpha_1=\alpha$, \emph{with no projection}.

\subsection{A conditional guarantee at a fixed level}

\begin{assumption}[Conditional independence]\label{ass:ci}
Given the features $E_{1:t}$, the scores $R_1,\dots,R_t$ are independent and
$R_s\mid E_{1:t}\sim\kappa(E_s,\cdot)$ for a fixed Markov kernel $\kappa$ (the law of the risk of a prompt with feature $e$).
Write $L_s=\kappa(E_s,\cdot)$.
\end{assumption}

Assumption~\ref{ass:ci} holds, for example, when prompts are independent draws and the model's
randomness is independent across prompts. It does not require the laws $L_s$ to be equal.

\begin{theorem}[Conditional coverage at a fixed level]\label{thm:cond}
Fix $t$ and $\alpha\in(0,1)$ and let $q_t^{+}$ be the threshold above with $\alpha_t\equiv\alpha$.
Under Assumption~\ref{ass:ci},
\[
 \mathbb P\bigl(R_t>q_t^{+}\ \big|\ E_{1:t}\bigr)\ \le\ \alpha+\frac{2}{1+m_t}\sum_{s<t}p_{t,s}\,
 d_{\mathrm{TV}}(L_s,L_t)\ \le\ \alpha+2\sum_{s<t}p_{t,s}\,d_{\mathrm{TV}}(L_s,L_t).
\]
In particular, if $L_s=L_t$ for every $s$ with $p_{t,s}>0$, then $\mathbb P(R_t>q_t^{+}\mid E_{1:t})\le\alpha$.
\end{theorem}

\begin{proof}
Condition on $E_{1:t}=e$; the weights $p_{t,s}$ are then fixed numbers. Apply Theorem~2 of
Barber, Cand\`es, Ramdas and Tibshirani (2023) and the remark following it (their Section~4.1) that the same
bound holds for nonexchangeable split conformal prediction, with the scores $R_s$ in place of the residuals and
no fitted function. The sequence is $(R_1,\dots,R_t)$ with conditional joint law $\mathcal L(R_{1:t}\mid e)$;
the fixed weights are $\hat w_s=p_{t,s}/m_t\in[0,1]$ for $s<t$ (their Section~3.1) and the test weight is $1$. Their normalised weights (their equation~(10)) are
$\hat w_s/(\sum_{j<t}\hat w_j+1)=p_{t,s}/(1+m_t)=\tilde w_{t,s}$, and their threshold (their equation~(11)) is the quantile
$q_t^{+}$ defined above. The theorem gives
$\mathbb P(R_t\le q_t^{+}\mid e)\ge 1-\alpha-\sum_{s<t}\tilde w_{t,s}\,d_{\mathrm{TV}}(\mathbf R,\mathbf R^{s})$,
where $\mathbf R=(R_1,\dots,R_t)$ and $\mathbf R^{s}$ is $\mathbf R$ with entries $s$ and $t$ swapped.
Under Assumption~\ref{ass:ci}, $\mathbf R$ has product law $\bigotimes_j L_j$ and $\mathbf R^{s}$ has the product law with
$L_s$ and $L_t$ exchanged, and these differ only in coordinates $s$ and $t$, so
$d_{\mathrm{TV}}(\mathbf R,\mathbf R^{s})\le d_{\mathrm{TV}}(L_s\otimes L_t,\ L_t\otimes L_s)\le 2\,d_{\mathrm{TV}}(L_s,L_t)$
(a product of two laws changes by at most the sum of the coordinate changes, by the triangle inequality).
Substituting gives the first bound; the second uses $1/(1+m_t)\le 1$.
\end{proof}

Theorem~\ref{thm:cond} holds at a fixed level and does not apply to the data-dependent level $\alpha_t$ of the online rule;
Theorem~\ref{thm:longrun} below covers that rule for every choice of weights, including weights that depend on past scores.

\begin{lemma}[Size of the correction]\label{lem:m}
$1/n_{\mathrm{eff},t}\le m_t\le n_{\mathrm{eff},t}^{-1/2}$.
\end{lemma}
\begin{proof}
$\sum_s p_{t,s}^2\le m_t\sum_s p_{t,s}=m_t$ gives the first inequality, and $m_t^2\le\sum_s p_{t,s}^2$ gives the second.
\end{proof}
The correction factor in $\ell_t$ therefore lies between $1+1/n_{\mathrm{eff},t}$ and $1+n_{\mathrm{eff},t}^{-1/2}$: concentrated weights are
charged in the level, not only in the variance of the weighted quantile.

\begin{corollary}[Drift controls the penalty]\label{cor:drift}
Suppose, in addition, $d_{\mathrm{TV}}\bigl(\mathcal L(R\mid E=e),\mathcal L(R\mid E=e')\bigr)\le \Lambda\|e-e'\|$ for all
$e,e'$ in the support of the features. Then, for the final (post-gate) weights,
\[
 \mathbb P(R_t>q_t^{+}\mid E_{1:t})\le\alpha+\frac{2\Lambda}{1+m_t}\,D_t,\qquad
 D_t=\sum_{s<t}p_{t,s}\,\|E_s-E_t\|.
\]
\end{corollary}

The quantity $D_t$ is the \emph{prompt drift} of the paper, computed with the weights that are actually used
to form the threshold (after the gate and recency mixing), so the diagnostic is the one that appears in the bound.
The Lipschitz constant $\Lambda$ is a property of the data-generating process; it is not estimated by the method,
and the bound is a statement about how the penalty scales, not a computable certificate.


Corollary~\ref{cor:drift} and Lemma~\ref{lem:m} show the trade-off that any drift-aware rule faces: concentrating weight on prompts close to
$E_t$ lowers $D_t$ and raises $m_t$ (lowers $n_{\mathrm{eff},t}$). Appendix~\ref{app:ess} gives the sampling error of the weighted quantile in the
no-drift case, which grows as $n_{\mathrm{eff},t}^{-1/2}$. Neither bound orders the rules, so the comparison between rules is empirical
(Section~\ref{sec:results}).

\subsection{A long-run guarantee for the online update}

\begin{theorem}[Long-run control, any weights]\label{thm:longrun}
For every $T\ge1$, every score sequence $R_1,\dots,R_T\in[0,1]$ and every choice of weights,
\[
 \Bigl|\frac1T\sum_{t=1}^T\mathrm{err}_t-\alpha\Bigr|\ \le\ \frac{\max\{\alpha_1,1-\alpha_1\}+\gamma}{T\gamma}.
\]
\end{theorem}

\begin{proof}
We first show $\alpha_t\in[-\gamma,1+\gamma]$ for all $t$, by induction from $\alpha_1=\alpha\in(0,1)$.
Suppose $\alpha_t\in[-\gamma,1+\gamma]$.
If $\alpha_t<0$ then $\ell_t=(1-\alpha_t)(1+m_t)>1$, so $q_t=+\infty$, $\mathrm{err}_t=0$ and
$\alpha_{t+1}=\alpha_t+\gamma\alpha\in(\alpha_t,\gamma\alpha)\subset[-\gamma,1+\gamma]$.
If $\alpha_t>1$ then $\ell_t<0$, so $q_t=-\infty$, $\mathrm{err}_t=1$ and
$\alpha_{t+1}=\alpha_t-\gamma(1-\alpha)\in[1-\gamma(1-\alpha),\alpha_t)\subset[-\gamma,1+\gamma]$.
If $\alpha_t\in[0,1]$ then $\alpha_{t+1}\in[-\gamma(1-\alpha),\,1+\gamma\alpha]\subset[-\gamma,1+\gamma]$.
Summing the update, $\alpha_{T+1}-\alpha_1=\gamma\sum_{t\le T}(\alpha-\mathrm{err}_t)$, so
$\frac1T\sum_t\mathrm{err}_t-\alpha=(\alpha_1-\alpha_{T+1})/(T\gamma)$, and
$\alpha_1-\alpha_{T+1}\le\alpha_1+\gamma$ and $\alpha_{T+1}-\alpha_1\le 1+\gamma-\alpha_1$, so
$|\alpha_1-\alpha_{T+1}|\le\max\{\alpha_1,1-\alpha_1\}+\gamma$, which gives the bound.
\end{proof}

This is Proposition~4.1 of Gibbs and Cand\`es (2021) specialised to our update; we give the proof because the
two conventions at the extremes ($\ell_t>1$ and $\ell_t<0$) are what the proof uses, and the correction factor
$1+m_t$ leaves them unchanged. The bound is the same for every weighting scheme: \emph{long-run marginal
coverage is a property of the update, not of PromptShift-CRC's weights}. The bound is $0.018$ to $0.024$ for the synthetic streams, $0.020$ for MMLU, $0.012$ for RealToxicityPrompts and $0.043$ for TruthfulQA.

\begin{remark}[What gating can and cannot do]\label{rem:gate}
Theorem~\ref{thm:longrun} gives every weighting scheme the same long-run bound, so long-run marginal coverage cannot distinguish schemes.
By Theorem~\ref{thm:cond} and Lemma~\ref{lem:m}, a gate that moves weight toward recent or similar prompts can reduce the drift penalty
$D_t$ while it reduces $n_{\mathrm{eff},t}$. We therefore do not claim that the gate or the representation weights improve
risk control; whether they improve local validity or sharpness is measured, not assumed.
\end{remark}
